\documentclass[a4paper,12pt]{article}
\usepackage[T2A]{fontenc}
\usepackage[utf8]{inputenc}
\usepackage[russian]{babel}
\usepackage{amsmath,amssymb,mathtools}
\renewcommand{\familydefault}{\sfdefault}
\usepackage{icomma}
\usepackage[margin=21mm,top=20mm,bottom=20mm,headheight=13pt]{geometry}
\usepackage{microtype}
\usepackage{tikz}
\usetikzlibrary{arrows.meta,calc,positioning,shapes.geometric}
\usepackage{tkz-euclide}
\usepackage{tabularx,booktabs,longtable}
\usepackage{enumitem}
\usepackage{hyperref}
\usepackage{parskip}
\usepackage{fancyhdr}
\usepackage{setspace}
\usepackage{xcolor}
\definecolor{accent}{HTML}{16676B}
\definecolor{darkaccent}{HTML}{164B4D}
\definecolor{textgray}{HTML}{363F41}
\definecolor{softgray}{HTML}{718081}
\definecolor{paleline}{HTML}{C2CFCD}
\hypersetup{colorlinks=true,linkcolor=darkaccent,urlcolor=accent,pdfauthor={Лёвин Артём Александрович},pdftitle={Чек-лист. Основы тригонометрии. Марк Гукин. 25.10.2026}}
\setlength{\parindent}{0pt}
\setlength{\parskip}{5pt}
\onehalfspacing
\setlist[itemize]{leftmargin=*,itemsep=2pt,topsep=3pt}
\setlist[enumerate]{leftmargin=*,itemsep=6pt,topsep=3pt}
\pagestyle{fancy}
\fancyhf{}
\fancyhead[L]{\footnotesize\color{softgray} ТРИГОНОМЕТРИЯ / I}
\fancyhead[R]{\footnotesize\color{softgray} 25.10.2026}
\fancyfoot[L]{\footnotesize\color{softgray} А. А. Лёвин \quad·\quad Марк Гукин}
\fancyfoot[R]{\footnotesize\color{softgray}\thepage}
\renewcommand{\headrulewidth}{0.3pt}
\renewcommand{\footrulewidth}{0pt}
\newcommand{\key}[1]{\textcolor{darkaccent}{\textbf{#1}}}
\newcommand{\dd}{^{\circ}}
\begin{document}
\begin{titlepage}
\thispagestyle{empty}
\begin{center}
\vspace*{32mm}
{\large\color{softgray}\textbf{МАТЕМАТИКА}\quad / \quad ИНДИВИДУАЛЬНОЕ ПОСОБИЕ}\par
\vspace{26mm}
{\color{darkaccent}\fontsize{35}{42}\selectfont\textbf{Чек-лист}}\par
\vspace{10mm}
{\fontsize{21}{27}\selectfont Основы тригонометрии}\par
\vspace{4mm}
{\large Единичная окружность, формулы, уравнения}\par
\vspace{26mm}
\begin{tikzpicture}
 \draw[paleline,line width=0.55pt] (-5.9,0)--(5.9,0);
 \fill[accent] (0,0) circle(2.5pt);
\end{tikzpicture}\par
\vspace{14mm}
{\large\textbf{Марк Гукин}}\par
\vspace{8mm}
{\large 25 октября 2026 года}\par
\vfill
{\small\color{softgray} Лёвин Артём Александрович}\par
\vspace{2mm}
{\small\color{textgray} эксклюзивно для Марка Гукина}\par
\vspace{8mm}
\begin{tikzpicture}[x=1cm,y=1cm]
\draw[accent,opacity=.65,line width=1pt] (-6,0) .. controls (-3.5,0.8) and (-1.6,-.45) .. (0,0)
 .. controls (2.2,.62) and (4,-.75) .. (6,.20);
\end{tikzpicture}
\end{center}
\end{titlepage}

\section*{01. Единичная окружность}
\key{Единичная окружность} имеет центр $O(0;0)$ и радиус $1$. Отправная точка~--- $(1;0)$: отсюда отсчитывают углы. Положительное направление~--- против часовой стрелки, отрицательное~--- по часовой стрелке.

\begin{center}
\begin{tikzpicture}[scale=2.0,>=Stealth]
  \tkzDefPoint(0,0){O}
  \tkzDefPoint(1,0){A}
  \tkzDefPoint(0,1){B}
  \tkzDefPoint(-1,0){C}
  \tkzDefPoint(0,-1){D}
  \tkzDefPoint(0.5,0.8660254){P}
  \draw[->,softgray,thin] (-1.3,0) -- (1.35,0) node[below] {$x$};
  \draw[->,softgray,thin] (0,-1.22) -- (0,1.28) node[left] {$y$};
  \tkzDrawCircle[draw=darkaccent,line width=.75pt](O,A)
  \draw[accent,line width=1pt] (O) -- (P);
  \draw[accent,densely dashed] (P) -- (0.5,0) (P) -- (0,0.8660254);
  \draw[->,accent,line width=.75pt] (0.35,0) arc[start angle=0,end angle=60,radius=.35];
  \node[accent,right] at (.39,.19) {$\alpha$};
  \tkzDrawPoints[fill=accent,size=3](O,A,B,C,D,P)
  \node[below left] at (O) {$O$};
  \node[below right,fill=white,inner sep=1pt] at (A) {$(1;0)$};
  \node[above left,fill=white,inner sep=1pt] at (B) {$(0;1)$};
  \node[below left,fill=white,inner sep=1pt] at (C) {$(-1;0)$};
  \node[below left,fill=white,inner sep=1pt] at (D) {$(0;-1)$};
  \node[above right,text=darkaccent,fill=white,inner sep=2pt] at (P) {$P\!\left(\tfrac12;\tfrac{\sqrt3}{2}\right)$};
\end{tikzpicture}
\end{center}
\[
\boxed{\text{угол }\alpha\ \longrightarrow\ P(x;y)\ \longrightarrow\ \cos\alpha=x,\quad\sin\alpha=y}
\]
Например, для $\alpha=60\dd$ получаем $P\left(\frac12;\frac{\sqrt3}{2}\right)$, поэтому $\cos60\dd=\frac12$, $\sin60\dd=\frac{\sqrt3}{2}$.

\key{Полный оборот:} $360\dd=2\pi$ радиан. Углы $-90\dd$ и $270\dd$ приводят в одну точку $(0;-1)$. Любые углы $\alpha$ и $\alpha+2\pi n$, где $n\in\mathbb Z$, имеют одинаковые синус и косинус.

\section*{02. Градусы и радианы}
\[\pi\ \text{рад}=180\dd,\qquad \alpha_{\text{рад}}=\alpha_{\text{град}}\cdot\frac{\pi}{180}.\]
\begin{center}
\renewcommand{\arraystretch}{1.30}
\begin{tabular}{@{}lccccc@{}}\toprule
Градусы & $0\dd$ & $30\dd$ & $45\dd$ & $60\dd$ & $90\dd$\\\midrule
Радианы & $0$ & $\pi/6$ & $\pi/4$ & $\pi/3$ & $\pi/2$\\\bottomrule
\end{tabular}
\end{center}

\clearpage
\section*{03. Точные значения синуса и косинуса}
Для углов $0\dd$, $30\dd$, $45\dd$, $60\dd$, $90\dd$ значения можно восстановить без заучивания: для синуса под корнем последовательно стоят $0,1,2,3,4$; для косинуса~--- в обратном порядке. Знаменатель всегда $2$.
\begin{center}
\renewcommand{\arraystretch}{1.80}
\begin{tabular}{@{}lccccc@{}}\toprule
$\alpha$ & $0$ & $\pi/6$ & $\pi/4$ & $\pi/3$ & $\pi/2$ \\ \midrule
$\sin\alpha$ & $0$ & $\frac12$ & $\frac{\sqrt2}{2}$ & $\frac{\sqrt3}{2}$ & $1$\\
$\cos\alpha$ & $1$ & $\frac{\sqrt3}{2}$ & $\frac{\sqrt2}{2}$ & $\frac12$ & $0$\\\bottomrule
\end{tabular}
\end{center}
\key{Углы, кратные $\pi$.} Для любого целого $m$:
\[
\sin(m\pi)=0,\qquad
\cos(m\pi)=(-1)^m=\begin{cases}1,&m\ \text{чётное},\\-1,&m\ \text{нечётное}.\end{cases}
\]
Например, $\sin(101\pi)=0$, $\cos(101\pi)=-1$, $\cos(-202\pi)=1$.

\section*{04. Простейшие тригонометрические уравнения}
Решить уравнение $\sin\alpha=c$ означает найти \emph{все} точки окружности с ординатой $c$. Для $\cos\alpha=c$ ищем точки с абсциссой $c$. Если $|c|>1$, решений нет.
\begin{center}
\renewcommand{\arraystretch}{1.42}
\begin{tabularx}{\linewidth}{@{}Xl@{}}\toprule
Уравнение & Все решения, $n\in\mathbb Z$ \\ \midrule
$\sin\alpha=1$ & $\displaystyle\alpha=\frac\pi2+2\pi n$ \\
$\sin\alpha=-1$ & $\displaystyle\alpha=-\frac\pi2+2\pi n$ \\
$\cos\alpha=1$ & $\alpha=2\pi n$ \\
$\cos\alpha=-1$ & $\alpha=\pi+2\pi n$ \\
$\sin\alpha=0$ & $\alpha=\pi n$ \\
$\cos\alpha=0$ & $\displaystyle\alpha=\frac\pi2+\pi n$ \\
\bottomrule\end{tabularx}
\end{center}
\key{Почему появляется $\pi n$?} Для $\sin\alpha=0$ подходят две противоположные точки: $0$ и $\pi$ (по модулю $2\pi$). Обе серии объединяются формулой $\alpha=\pi n$.

\key{Проверка:} ответ должен учитывать все подходящие точки и все полные обороты. Не смешивайте градусы и радианы в одной записи.

\clearpage
\section*{05. Основное тригонометрическое тождество}
\[
\boxed{\sin^2\alpha+\cos^2\alpha=1}\qquad \text{при любом действительном }\alpha.
\]
\key{Приём 1.} Если коэффициенты одинаковы, вынесите общий множитель:
\[
5\sin^2\alpha+5\cos^2\alpha
=5(\sin^2\alpha+\cos^2\alpha)=\boxed{5}.
\]
\key{Приём 2.} Если коэффициенты различны, выделите полную сумму квадратов:
\begin{align*}
5\sin^2\alpha+7\cos^2\alpha
&=5(\sin^2\alpha+\cos^2\alpha)+2\cos^2\alpha\\
&=\boxed{5+2\cos^2\alpha}.
\end{align*}
\key{Приём 3.} Подставляйте тождество в обратном направлении:
\[
\boxed{8=8\sin^2 50\dd+8\cos^2 50\dd}.
\]
Углы внутри двух квадратов должны совпадать. В тождестве стоят именно \emph{квадраты}, а между ними~--- \emph{плюс}.

\section*{06. Синус и косинус суммы и разности}
\begin{align*}
\sin(\alpha\pm\beta)&=\sin\alpha\cos\beta\ \pm\ \cos\alpha\sin\beta,\\
\cos(\alpha\pm\beta)&=\cos\alpha\cos\beta\ \mp\ \sin\alpha\sin\beta.
\end{align*}
\key{Особенность:} в формуле косинуса знак между произведениями противоположен знаку между углами.

\key{Пример 1.} Вычислим $\sin75\dd$, представив $75\dd=30\dd+45\dd$:
\begin{align*}
\sin75\dd
 &=\sin30\dd\cos45\dd+\cos30\dd\sin45\dd\\
 &=\frac12\cdot\frac{\sqrt2}{2}+\frac{\sqrt3}{2}\cdot\frac{\sqrt2}{2}
 =\boxed{\frac{\sqrt2+\sqrt6}{4}}.
\end{align*}
\key{Пример 2.} Вычислим $\cos135\dd$, представив $135\dd=90\dd+45\dd$:
\begin{align*}
\cos135\dd
 &=\cos90\dd\cos45\dd-\sin90\dd\sin45\dd\\
 &=0\cdot\frac{\sqrt2}{2}-1\cdot\frac{\sqrt2}{2}
 =\boxed{-\frac{\sqrt2}{2}}.
\end{align*}
\key{Самопроверка.} $\sin15\dd>0$, а $\cos135\dd<0$ по расположению соответствующих точек на окружности. Знак точного ответа должен этому соответствовать.

\clearpage
\section*{Алгоритм. Решение простейшего тригонометрического уравнения}
\vspace{8mm}
\begin{center}
\begin{tikzpicture}[
 >=Stealth,
 flow/.style={draw=darkaccent,rounded corners=3pt,align=center,text width=9.1cm,inner sep=8pt,minimum height=14mm},
 condition/.style={draw=darkaccent,diamond,aspect=2.7,align=center,text width=3.5cm,inner sep=3pt},
 terminal/.style={draw=darkaccent,ellipse,align=center,text width=3.4cm,inner sep=6pt},
 arr/.style={->,line width=.75pt,draw=accent},
 node distance=10mm
]
\node[terminal] (start) {Уравнение $\sin\alpha=c$ или $\cos\alpha=c$};
\node[condition,below=12mm of start] (check) {$|c|\le1$?};
\node[terminal,left=36mm of check] (none) {Решений нет};
\node[flow,below=12mm of check] (axis) {На единичной окружности проведите линию $y=c$ для синуса или $x=c$ для косинуса};
\node[flow,below=11mm of axis] (points) {Найдите все точки пересечения с окружностью и определите соответствующие углы};
\node[flow,below=11mm of points] (cycles) {К каждому базовому углу добавьте $2\pi n$, где $n\in\mathbb Z$};
\node[flow,below=11mm of cycles] (combine) {Если возможно, объедините серии: $\sin\alpha=0\Rightarrow\alpha=\pi n$; $\cos\alpha=0\Rightarrow\alpha=\frac\pi2+\pi n$};
\node[terminal,below=12mm of combine] (end) {Запишите все решения};
\draw[arr] (start)--(check);
\draw[arr] (check.west)--node[above,softgray]{нет}(none.east);
\draw[arr] (check.south)--node[right,softgray]{да}(axis.north);
\draw[arr] (axis)--(points);
\draw[arr] (points)--(cycles);
\draw[arr] (cycles)--(combine);
\draw[arr] (combine)--(end);
\end{tikzpicture}
\end{center}
\clearpage

\section*{Тренировочный блок / Путь А}
\textcolor{softgray}{10 заданий. Выполняйте по порядку. В заданиях с уравнениями указывайте все решения и $n\in\mathbb Z$.}
\begin{enumerate}[label=\textcolor{darkaccent}{\bfseries\arabic*.},leftmargin=12mm,itemsep=11pt]
\item На одной единичной окружности отметьте точки для углов $45\dd$, $-45\dd$, $180\dd$, $210\dd$, $270\dd$. Запишите точные координаты каждой точки.
\item Переведите в радианы углы $30\dd$, $135\dd$, $225\dd$, $-90\dd$. Переведите в градусы углы $\frac\pi6$ и $\frac{5\pi}{4}$.
\item Вычислите по единичной окружности:
\[\sin\frac\pi2;\qquad \cos\pi;\qquad \sin\left(-\frac\pi2\right);\qquad \cos2\pi.\]
\item Восстановите табличные значения:
\[\sin\frac\pi6;\qquad \cos\frac\pi3;\qquad \sin\frac\pi4;\qquad \cos\frac\pi6.\]
\item Не выполняя многократных оборотов, найдите:
\[\sin(101\pi);\qquad \cos(-202\pi);\qquad \cos(103\pi).\]
\item Решите уравнения: \quad а) $\sin\alpha=-1$; \quad б) $\cos\alpha=1$.
\item Решите уравнения: \quad а) $\sin\alpha=0$; \quad б) $\cos\alpha=0$.
\item Упростите: \quad а) $7\sin^2\alpha+7\cos^2\alpha$;\\
\hspace*{20mm}б) $5\sin^2\alpha+7\cos^2\alpha$.
\item Представьте число $9$ как сумму выражений с $\sin^2 50\dd$ и $\cos^2 50\dd$. Затем упростите $3-3\cos^2\alpha$.
\item Используя формулы суммы и разности, вычислите точно:
\[\sin15\dd;\qquad \cos105\dd.\]
\end{enumerate}

\clearpage
\section*{Ответы к тренировочному блоку}
\textcolor{softgray}{Проверьте ответы после самостоятельного решения. Промежуточные преобразования сверяйте с примерами выше.}
\renewcommand{\arraystretch}{1.42}
\begin{longtable}{@{}p{9mm}p{\dimexpr\linewidth-15mm\relax}@{}}
\toprule
№ & Ответ\\\midrule\endfirsthead
\toprule № & Ответ (продолжение)\\\midrule\endhead
1 & $45\dd:\left(\frac{\sqrt2}{2};\frac{\sqrt2}{2}\right)$; $-45\dd:\left(\frac{\sqrt2}{2};-\frac{\sqrt2}{2}\right)$;\\
  & $180\dd:(-1;0)$; $210\dd:\left(-\frac{\sqrt3}{2};-\frac12\right)$; $270\dd:(0;-1)$.\\[3pt]
2 & $\frac\pi6;\ \frac{3\pi}{4};\ \frac{5\pi}{4};\ -\frac\pi2;\quad 30\dd;\ 225\dd$.\\[3pt]
3 & $1;\ -1;\ -1;\ 1$.\\[3pt]
4 & $\frac12;\ \frac12;\ \frac{\sqrt2}{2};\ \frac{\sqrt3}{2}$.\\[3pt]
5 & $0;\ 1;\ -1$.\\[3pt]
6 & а) $\displaystyle\alpha=-\frac\pi2+2\pi n$; \quad б) $\alpha=2\pi n$, $n\in\mathbb Z$.\\[3pt]
7 & а) $\alpha=\pi n$; \quad б) $\displaystyle\alpha=\frac\pi2+\pi n$, $n\in\mathbb Z$.\\[3pt]
8 & а) $7$; \quad б) $5+2\cos^2\alpha$.\\[3pt]
9 & $9\sin^2 50\dd+9\cos^2 50\dd$; \quad $3\sin^2\alpha$.\\[3pt]
10 & $\displaystyle\sin15\dd=\frac{\sqrt6-\sqrt2}{4}$; \quad $\displaystyle\cos105\dd=\frac{\sqrt2-\sqrt6}{4}$.\\
\bottomrule
\end{longtable}
\vspace{9mm}
\textcolor{darkaccent}{\textbf{Контроль освоения:}}
\begin{itemize}[label=$\square$,leftmargin=*,itemsep=1pt]
\item Я различаю $x=\cos\alpha$ и $y=\sin\alpha$.
\item Я перевожу основные углы в радианы и обратно.
\item Я восстанавливаю таблицу значений без подсказки.
\item Я записываю все корни простейших уравнений.
\item Я применяю $\sin^2\alpha+\cos^2\alpha=1$ в обе стороны.
\item Я верно выбираю знаки в формулах суммы и разности.
\end{itemize}
\end{document}
